\documentclass[11pt,a4paper]{article}
\usepackage[margin=24mm]{geometry}
\usepackage{amsmath,amssymb,booktabs,microtype}
\usepackage[hidelinks]{hyperref}
\newcommand{\dd}{\mathrm d}
\title{Route G2-P: Physical Anchors and a Visible Probe Candidate}
\author{Scale identification, a conditional Higgs coupling, and calibration inputs}
\date{23 September 2026}
\begin{document}
\maketitle
\begin{abstract}
The existing readout forecast cannot select an absolute energy scale.
We derive a spectral anchor test and a symmetry-compatible visible-sector
probe, retaining all missing experimental inputs explicitly. The probe
is a candidate four-dimensional effective interaction, not an activated
hardware design, a Standard Model particle identification, or a Route-F
portal. Existing G2 results remain the zero-probe regression baseline.
\end{abstract}

\section{What the dimensionless forecast cannot determine}
Use $\hbar=c=1$, $\Lambda=R^{-1}$, and hats for the frozen $R=1$ card.
The existing engineering ratios specify
\begin{align}
 (M_5,M_D,p_0,\Omega,\sigma_i,\sigma_d,c)
 &=\Lambda(.5,5,.2,.2,.1,.05,1.25892910),\\
 \alpha&=.25,\quad g=.05,\quad\epsilon=.5,\quad
 \kappa_5=2\pi g/\Lambda .
\end{align}
These are relations, not a choice of eV or GeV. They give
$m_n=\Lambda\sqrt{.25+(n+.25)^2}$ and
$\mu_j=\Lambda\sqrt{25+j^2}$. Restore physical lengths and times as
$\hbar c/\Lambda$ and $\hbar/\Lambda$ if conventional units are used.
For a long-time sideband $q$, the unchanged tree kernel is
\begin{equation}
 K_{lmq}=\frac{|g_q|^2p_f}{16\pi E_nE_l(E_n+E_l+q\Omega)}
          =\Lambda^{-2}\widehat K_{lmq},
 \qquad g_{\pm1}=g\epsilon/2,\quad g_0=g.
\end{equation}
Thus $K$ and cross sections scale as $\Lambda^{-2}$; mean work per
scattered event scales as $\Lambda$, and $\sum w_l K_{lmq}q\Omega$
as $\Lambda^{-1}$. The last quantity is not power without luminosity.

In the rate-weighted beam mixture, $f_\Lambda(p)=\Lambda^{-1}\hat f(p/\Lambda)$
and ${\cal R}_\Lambda(x|u)=\Lambda^{-1}\hat{\cal R}(x/\Lambda|u/\Lambda)$.
Consequently $h_{s,\Lambda}(x)=\Lambda^{-3}\hat h_s(x/\Lambda)$,
$Z_\Lambda=\Lambda^{-2}\hat Z$, and
\begin{equation}
 \frac{\int\min(h_{+,\Lambda},h_{-,\Lambda})\dd x}{Z_\Lambda}
 =\frac{\int\min(\hat h_+,\hat h_-)\dd\hat x}{\hat Z}
 =.022022469\ldots .
\end{equation}
Scaling the frozen classifier and its offset/noise tolerances gives the
same result. This is a scale-identifiability limitation of these inputs,
not a symmetry of the full SM with its measured masses held fixed.
Merely specifying a pump frequency could set a scale only after an
actual frequency and the imposed $\Omega/\Lambda=.2$ relation are supplied;
that is an external anchor, not a predicted radius.

\section{A correlated spectrum, not one fitted scale per particle}
For calibrated masses at assigned consecutive \emph{signed} labels
$n=-1,0,1$, put $y_n=m_n^2$ and fit $y_n=An^2+Bn+C$:
\begin{equation}
 A=\frac{y_1+y_{-1}-2y_0}{2},\qquad
 B=\frac{y_1-y_{-1}}{2},\qquad C=y_0.
\end{equation}
Then $\Lambda=\sqrt A$, $\alpha=B/(2A)$, and
$M_{\rm eff}^2=C-B^2/(4A)$. The original nonnegative-bulk-mass
model requires $A>0$, $M_{\rm eff}^2\geq0$. Large gauge transformations
relabel $\alpha\mapsto\alpha+k$, $n\mapsto n-k$; sorted masses are
not automatically adjacent signed labels. Under the fully frozen ratios,
one externally identified mass suffices to set $\Lambda=m_n/\hat m_n$,
but does not test the other assumed ratios. Without freezing them, a
fourth signed mass supplies an out-of-fit condition:
\begin{equation}\label{eq:closure}
 r=y_2-3y_1+3y_0-y_{-1}=0,\qquad
 \operatorname{Var}(r)=d^T V_y d,\quad d=(-1,3,-3,1)^T .
\end{equation}
$V_y$ is the calibrated squared-mass covariance; to first order
$V_y=\operatorname{diag}(2m)V_m\operatorname{diag}(2m)$.
An experimental rejection requires these errors and a declared theory
error, not floating-point equality. Loops can modify this tree identity.
The supplied numerical tests use synthetic anchors only, recover the
input scale, and reject a deliberately distorted fourth mass.

\section{A visible-sector coupling: declare the extension first}
Treat the towers as SM gauge singlets below an explicit EFT cutoff.
One circle-uniform candidate is
\begin{equation}\label{eq:portal}
 \Delta{\cal L}_4=-H^\dagger H
 \left(\lambda_\Phi\sum_n|\phi_n|^2+
       \lambda_X\sum_j|x_j|^2\right),\qquad \lambda_\Phi,\lambda_X\geq0.
\end{equation}
This preserves both original scalar charges, the external flat-$U(1)$
covariance, and compact momentum. It follows from
$-\int R\dd\theta\,H^\dagger H(\lambda_\Phi|\Phi|^2+\lambda_X|X|^2)$
with $H(x)$ independent of $\theta$, because normalized Fourier overlaps
give $\delta_{nm}$. As written it is local in four dimensions, with a
circle-smeared Higgs degree of freedom; it is \emph{not} a demonstrated
strictly five-dimensional local SM completion. A bulk Higgs zero mode
would instead give $\lambda_4=\lambda_5/(2\pi R)$ and require its other
modes, gauge sector and profiles to be specified. No such completion is
silently assumed. Scalar Higgs-density couplings are standard
\cite{cline}; the particular tower and readout here are our hypothesis.

With $H=(0,(v+h)/\sqrt2)^T$, the tree masses and vertices become
\begin{align}
 m_n^2&=M_{5,0}^2+\lambda_\Phi v^2/2+(n+\alpha)^2\Lambda^2,\\
 \mu_j^2&=M_{D,0}^2+\lambda_X v^2/2+j^2\Lambda^2,\\
 \Delta{\cal L}_4&\supset-\sum_{S,k}
   (\lambda_S v h+\lambda_S h^2/2)|S_k|^2.
\end{align}
The common shifts preserve \eqref{eq:closure}, but spectral intercepts
cannot separate a pre-EWSB mass from $\lambda_Sv^2/2$. The SM value of
$v$ does not fix $R$ when these parameters remain independent. Keeping
the old physical masses would require an explicit mass matching
condition, not an unnoticed subtraction. Nonnegative pre-EWSB masses
and couplings give a simple classically safe sufficient choice; negative
matched masses require a new vacuum analysis and are not accepted here.
In particular, retaining $M_{D,\mathrm{eff}}^2=25\Lambda^2$ gives
\begin{equation}\label{eq:stability}
 M_{D,0}^2=25\Lambda^2-\lambda_Xv^2/2\geq0
 \quad\Longrightarrow\quad \lambda_X\leq50\Lambda^2/v^2.
\end{equation}
This is substantive: with no $|X|^4$ in the retained polynomial action,
a negative $M_{D,0}^2$ runs away along $H=\Phi=0$ and a constant $X$.
Adding a stabilizing self-quartic would be a new parameter and vacuum
problem. The analogous nonnegative $M_{5,0}^2$ condition is sufficient,
not necessary, because a noninteger holonomy supplies a gradient gap.

If both couplings are nonzero, Higgs exchange also adds the diagonal
tree amplitude
\begin{equation}
 {\cal M}(\phi_n x_l\to\phi_n x_l)
 =-g+\frac{\lambda_\Phi\lambda_Xv^2}{m_h^2-t}
 \quad\hbox{(undriven elastic component)}.
\end{equation}
The zero-mode Higgs supplies no direct $n\ne m$ conversion, but elastic
rates and hence total branching fractions change. The minimal
detector-facing matching choice is $\lambda_X\ne0$, $\lambda_\Phi=0$
at a declared scale; mixed $g\lambda_X$ loops allow a $\lambda_\Phi$
counterterm, so this zero is not protected to all orders. Neither choice
has been activated in the existing numerical collision forecasts.

\section{From a coupling to measurable transfers}
For one complex mode $S$ of physical mass $m_S$, define
${\cal L}_{hNN}=-(f_Nm_N/v)h\bar NN$. At small nucleon momentum
transfer, treating $f_N$ as constant, a single-Higgs exchange gives
\begin{align}
 \overline{|{\cal M}|^2}
 &=\frac{\lambda_S^2 f_N^2m_N^2}{(m_h^2-t)^2}(4m_N^2-t),\\
 \frac{\dd\sigma}{\dd t}
 &=\frac{\overline{|{\cal M}|^2}}{16\pi{\cal K}(s,m_S^2,m_N^2)},
 \quad {\cal K}(a,b,c)=a^2+b^2+c^2-2ab-2ac-2bc .
\end{align}
The trace identity is
$\tfrac12\operatorname{Tr}[(\not k'+m_N)(\not k+m_N)]
=4m_N^2-t$. In the nonrelativistic contact limit, the allowed interval
has width $4p_*^2={\cal K}/s$, yielding
\begin{equation}\label{eq:sigma}
 \sigma_{SN}^{\rm NR}=
 \frac{\lambda_S^2 f_N^2m_N^4}{4\pi m_h^4(m_S+m_N)^2}.
\end{equation}
This agrees with the standard scalar normalization \cite{cline}.
It is a cross section per specified incident mode, not a cosmological
dark-matter event rate or a momentum resolution. Nuclear form factors,
coherence, target excitations and nonconstant hadronic response have
not been replaced by a nucleon formula.

For a stationary free nucleon and incident momentum $p$,
$s=m_S^2+m_N^2+2m_N\sqrt{m_S^2+p^2}$,
$t=-2m_NE_R$ and
\begin{equation}
 \frac{\dd\sigma}{\dd E_R}=2m_N\frac{\dd\sigma}{\dd t},\qquad
 0\leq E_R\leq\frac{{\cal K}(s,m_S^2,m_N^2)}{2m_Ns}.
\end{equation}
A single recoil deposit does not uniquely determine incident $p$;
the scattering angle is unobserved. A real response must integrate this
kernel with target response, efficiency and geometry. Even a nonzero
coupling therefore does not establish the old Gaussian $\sigma_d=.05\Lambda$.
For dilute independent targets, column density ${\cal C}=n_Td$ gives
the ideal first-interaction probability $1-e^{-{\cal C}\sigma}$.
This is only an opacity estimate, not accepted-event efficiency,
spectrometer performance, or a multiple-scattering transport model.

The same coupling yields an independent consistency test. For a complex
tower with all kinematically open modes inside EFT control,
\begin{equation}
 \Gamma_{h\to S\bar S}=\frac{\lambda_S^2v^2}{16\pi m_h}
 B_S,\qquad B_S=\sum_{k:2m_{S,k}<m_h}
 \sqrt{1-4m_{S,k}^2/m_h^2}.
\end{equation}
Each signed mode is a separate complex field; particle/antiparticle
are already counted once in its distinct two-body final state. There
is no identical-real-scalar factor $1/2$. An invisible-width limit is
applicable only if these states actually escape the relevant detector.
If $B_S>0$ and an applicable width budget is $\Gamma_{\max}$, eliminate
the same $\lambda_S$ between decay and scattering to obtain
\begin{equation}\label{eq:bound}
 \sigma_{SN}^{\rm NR}\leq
 \frac{4 f_N^2m_N^4\Gamma_{\max}}
 {v^2m_h^3(m_S+m_N)^2B_S}.
\end{equation}
Other open invisible sectors consume the budget and tighten this
necessary bound. If $B_S=0$, this decay gives no bound, not infinite
evidence for detectability. ATLAS provides a real example of a Higgs
invisible search \cite{atlas}; no published single-species exclusion
is applied without checking the tower, lifetime, production and
acceptance assumptions. No local dark-matter abundance is assumed.
Equation~\eqref{eq:bound} is evaluated at specified physical masses.
When bare masses instead are fixed and $\lambda_S$ varied, those masses
and $B_S$ vary too; the bound is implicit. For the frozen physical
$X$ intercept one must also impose \eqref{eq:stability}. Detectability
cannot be rescued by arbitrarily increasing $\lambda_X$ while ignoring
the vacuum or the summed tower width.

\section{Alternatives, calibration, and a finite next decision}
A brane-localized Higgs-density coupling is different physics:
\begin{equation}
 -\beta_S H^\dagger H |S(x,\theta_b)|^2
 =-H^\dagger H\sum_{m,n}\frac{\beta_S}{2\pi R}
 e^{i(n-m)\theta_b}S_m^*S_n,\qquad [\beta_S]=-1.
\end{equation}
It is off diagonal (rank one in a finite retained tower), mixes the
vacuum masses after EWSB, and lets the fixed brane absorb compact
momentum. The old selection rule then applies only to the isolated
conversion vertex, not the whole apparatus. It cannot be appended as
an innocent resolution model. A uniform density probe instead preserves
the $j\leftrightarrow-j$ mass and direct-response degeneracy of $X$;
sign inference still needs the collision correlations, not merely mass.

The accompanying calibration card requests one physical energy anchor
with units and covariance, a specified interaction/target, beam and
frame definitions, efficiencies, and a joint mass/momentum response.
For a finite event rate the measured target column or luminosity is
also necessary. A sideband tag requires an independent calibrated
response matrix, not the drive clock phase. Synthetic scale rescalings
and synthetic portal tests are explicitly separated from measured data.

The bounded next decision is physical: (i) choose this neutral
Higgs-coupled sector and an externally justified scale, then test
\eqref{eq:bound} against a specified target requirement; or (ii) supply
an existing apparatus response and compare it with the frozen contract;
or (iii) choose an analogue-mode platform, explicitly giving up a claim
of new elementary-particle masses. No additional precision scan is a
substitute. If a required scattering strength violates a justified
width/production constraint, reject that candidate rather than fit
each mode separately. A dimensional radius-stabilization theory is a
separate hypothesis, not needed to perform this empirical test.

\begin{thebibliography}{9}
\bibitem{cline} J.~M.~Cline, K.~Kainulainen, P.~Scott and C.~Weniger,
\emph{Update on scalar singlet dark matter}, Phys. Rev. D 88, 055025 (2013),
\href{https://arxiv.org/abs/1306.4710}{arXiv:1306.4710},
Eqs.~(3), (23), (26). Used for normalization, not old exclusion values.
\bibitem{atlas} ATLAS Collaboration, \emph{Combination of searches for
invisible decays of the Higgs boson using 139 fb$^{-1}$},
Phys. Lett. B 842, 137963 (2023),
\href{https://arxiv.org/abs/2301.10731}{arXiv:2301.10731}.
An example of a conditional experimental constraint, not calibration
data for the Route-G recoil field.
\end{thebibliography}
\end{document}
